\documentclass[10pt,a4paper]{amsart}
\usepackage[margin=19mm]{geometry}
\usepackage{amsmath,amssymb,mathtools}
\usepackage[scr=rsfs,cal=euler]{mathalfa}
\usepackage{xfrac}
\usepackage[unicode,hidelinks,bookmarks=false]{hyperref}
\newcommand{\ie}{i.e.\ }
\newcommand{\cf}{cf.\ }
\newcommand{\resp}{respectively\ }
\newcommand{\Subsection}{\subsection}
\providecommand{\nobreakdash}{\nobreak\hbox{-}}
\setlength{\emergencystretch}{2em}
\multlinegap0pt
\usepackage[shortlabels]{enumitem}
\usepackage{dsfont}
\usepackage{xcolor}
\usepackage{amssymb}
\usepackage{mathtools}
\usepackage{cancel}
\usepackage{float}
\usepackage{bm}
\usepackage{tikz}
\newtheorem{proposition}{Proposition}
\newcommand{\E}{\mathbb{E}}
\newcommand{\R}{\mathbb{R}} % real number
\newcommand{\cQ}{\mathcal{Q}}
\DeclareMathOperator{\cleq}{\,\preceq_{\rm c}\,}
\newcommand*{\rd}{\mathrm{d}}
\title{Bid--Ask Martingale Optimal Transport}
\author{Bryan Liang, Marcel Nutz, Shunan Sheng, Valentin Tissot-Daguette}
\date{Source: arXiv:2603.24605v2 (CC BY 4.0)}
\begin{document}
\maketitle

\noindent\emph{Source and licence.} Bid--Ask Martingale Optimal Transport. Version arXiv:2603.24605v2 (CC BY 4.0), distributed by arXiv under the Creative Commons Attribution 4.0 International licence, http://creativecommons.org/licenses/by/4.0/, as recorded in the arXiv licence metadata for this exact version and shown on the abstract page https://arxiv.org/abs/2603.24605v2. Full licence notice: \texttt{licenses/arxiv-2603.24605-CC-BY-4.0.txt}.\par\smallskip

\noindent\textit{Editorial scope.} Counted unit: the article's proposition prop:digital-call-spread (source0.tex lines 2194--2203). The article's complete original proof of it is delivered with it (source0.tex lines 2204--2280, one \texttt{proof} environment of 77 lines). No mathematics is written, repaired or re-derived: the statement and the proof are the article's own lines, in the article's own notation, and no retained line is substituted or re-flowed.  No further source-local context is needed: every cross-reference of every retained line resolves inside the two delivered spans.  Nothing was omitted from any retained span, so the package carries no omission record.  No retained line contains a citation, so the wrapper carries no bibliography and no bibliography entry.  No retained line contains a figure environment, an image-inclusion command or drawing code, so no image is delivered and the package declares no asset dependency.  Editorial formatting only, in addition: an AMS \texttt{amsart} preamble that restates the article's own declarations and macros used by the retained lines, namely the environments \texttt{proposition} on separate counters, together with the standard packages those lines need.  The retained environments print in this document as Proposition 1 in order of appearance, whereas the article prints the same items with the numbering of its own full document.  The complete native source of the article as distributed in the arXiv e-print for this exact version is bundled unchanged as \texttt{sources/197158/source0.tex}.
\begin{proposition}\label{prop:digital-call-spread}
Let $\mu^{\rm b}\cleq \mu^{\rm a}$ satisfy $c^{\rm a}(K,T)>c^{\rm b}(K,T)$. Let $h(x) = \mathds{1}_{\{x\geq K\}}$ be the payoff function.
It follows that the function $[0,K)\ni L\mapsto f_K(L)$ attains a minimizer $L^\star$. Then
\begin{enumerate}[label = (\roman*)]
    \item $\frac{1}{K-L^\star}((x-L^\star)^+, (x-K)^+) \in \Psi^{\rm b,a}(h)$ is a dual optimizer.
\item $\mu^\star = \mu^{\rm a}|_{[0,L^\star)} + (\mu^{\rm a}([L^\star, \infty)) - q^\star) \delta_{L^\star} + (q^\star - \mu^{\rm b}((K, \infty)) ) \delta_{K} +  \mu^{\rm b}|_{(K,\infty)}$ is a primal optimizer, where $q^\star =  f_K(L^\star)$.
\item $P(h)=D(h)= q^\star$.
\item Write $c^\star(k,T):=\E_{\mu^\star}[(X-k)^+]$, it holds that $c^\star(L^\star,T)=c^{\rm a}(L^\star,T)$, $c^\star(K,T)=c^{\rm b}(K,T)$.
\end{enumerate}
\end{proposition}

\begin{proof}
Since $c^{\rm a}(K,T)>c^{\rm b}(K,T)$, we know that $f_K(L)\to\infty$ as
$L\uparrow K$. Together with the continuity of call prices in strikes, this implies the existence of a minimizer $L^\star$. 

Set
$\ell(k):=c^{\rm b}(K,T)+q^\star(K-k)$. Minimality of $L^\star$ gives
$\ell\le c^{\rm a}(\cdot,T)$ on $[0,K]$, while convexity of
$c^{\rm b}(\cdot,T)$ gives $c^{\rm b}(\cdot,T)\le\ell$ on
$[L^\star,K]$. Hence
\[
c^\star(k,T):=
\begin{cases}
c^{\rm a}(k,T),&k\le L^\star,\\
\ell(k),&L^\star<k<K,\\
c^{\rm b}(k,T),&k\ge K,
\end{cases}
\]
is piecewise convex and satisfies
$c^{\rm b}\le c^\star\le c^{\rm a}$. To verify its global convexity, let $s=-q^\star$ denote the slope of
$\ell$. Since $\ell\le c^{\rm a}(\cdot,T)$ with equality at $L^\star$,
we have $\partial^-c^{\rm a}(L^\star,T)\le s$. Similarly, since
$c^{\rm b}(\cdot,T)\le\ell$ on $[L^\star,K]$ with equality at $K$, we
have $s\le\partial^-c^{\rm b}(K,T)\le\partial^+c^{\rm b}(K,T)$. Thus
the one-sided slopes of $c^\star$ are nondecreasing at $L^\star$ and
$K$, proving that $c^\star(\cdot,T)$ is convex.

Now we construct a primal optimizer from $c^\star$. In view of the identities
$\partial^-c^{\rm a}(L^\star,T)
=-\mu^{\rm a}([L^\star,\infty))$ and
$\partial^+c^{\rm b}(K,T)
=-\mu^{\rm b}((K,\infty))$, we have $\mu^{\rm a}([L^\star,\infty))-q^\star\ge0$ and
$q^\star-\mu^{\rm b}((K,\infty))\ge0$. If $L^\star=0$, the first
inequality follows instead from $q^\star = f_K(0)\le1$. We may
therefore define
\[
\mu^\star
=
\mu^{\rm a}\big|_{[0,L^\star)}
+\big(\mu^{\rm a}([L^\star,\infty))-q^\star\big)\delta_{L^\star}
+\big(q^\star-\mu^{\rm b}((K,\infty))\big)\delta_K
+\mu^{\rm b}\big|_{(K,\infty)}.
\]
It is clear that $\mu^\star$ is a probability
measure. For $L^\star\le k\le K$, direct computation gives
\[
\begin{aligned}
\E_{\mu^\star}[(X-k)^+]
&=\big(q^\star-\mu^{\rm b}((K,\infty))\big)(K-k)
  +\int_{(K,\infty)}(x-k)\,\mu^{\rm b}(\rd x)\\
&=c^{\rm b}(K,T)+q^\star(K-k)=\ell(k).
\end{aligned}
\]
For $k\ge K$, the same expectation is $c^{\rm b}(k,T)$. At
$k=L^\star$, the preceding computation gives
$\E_{\mu^\star}[(X-L^\star)^+]=c^{\rm a}(L^\star,T)$, while
$\mu^\star([L^\star,\infty))
=\mu^{\rm a}([L^\star,\infty))$. Splitting the call payoff at
$L^\star$ therefore gives
$\E_{\mu^\star}[(X-k)^+]=c^{\rm a}(k,T)$ for $k\le L^\star$.
Consequently,
$\E_{\mu^\star}[(X-k)^+]=c^\star(k,T)$ for all $k\in\R_+$.
Since $c^{\rm b}\le c^\star\le c^{\rm a}$, we conclude that
$\mu^\star\in\cQ(\mu^{\rm b},\mu^{\rm a})$. Moreover,
$\mu^\star([K,\infty))=q^\star$.
\\
Finally, let
$\psi^\star(x)=((x-L^\star)^+-(x-K)^+)/(K-L^\star)$. Then
$\psi^\star\ge h$ and its cost is $q^\star$. Together with weak
duality, this yields
\[
q^\star\le P(h)\le D(h)\le q^\star.
\]
Thus $\mu^\star$ and the stated call spread are primal and dual
optimizers, respectively. The identities
$c^\star(L^\star,T)=c^{\rm a}(L^\star,T)$ and
$c^\star(K,T)=c^{\rm b}(K,T)$ follow directly from the construction.
\end{proof}

\end{document}
